\section{Scope and Robustness}
\label{sec:secondary}

The residual-rent result is exact inside the canonical Cournot microfoundation, but R7 does not certify it as a general theorem across demand systems or coalition institutions. This section makes that boundary explicit. The negative results are not repairs to a failed proof; they identify which parts of the economic mechanism are portable and which parts are specification dependent.

\subsection{Zero-network boundary}
\label{subsec:zero-network}

At the boundary $v=0$, the pre-adoption member payoff difference simplifies to
\begin{equation}
    W_M^{SU,N}-W^{IS}
    =
    \frac{c(13c-6)}{32}.
    \label{eq:zero-network-gap}
\end{equation}
For every $0<c<1/3$ this expression is strictly negative. Thus the canonical political reversal cannot begin at zero network effects because the regional union is not initially preferred.

This does not mean that private one-way adoption disappears when $v=0$. The adoption thresholds reduce to
\begin{equation}
T_A(0)=\frac{3c(2-3c)}{16},
\qquad
T_U(0)=T_W(0)=\frac{3c(2-c)}{16},
\label{eq:zero-network-adoption-thresholds}
\end{equation}
and
\begin{equation}
T_U(0)<2T_A(0)
\qquad (0<c<1/3).
\end{equation}
Hence a positive fixed-cost interval can still generate outsider-only adoption. The boundary therefore separates the private-adoption logic from the initial government-ranking condition.

R4 gives the stronger fixed-$c$ characterization. Let
\[
\bar v(c)
=
\min\left\{\frac14,\frac{1-3c}{3(1-c)}\right\}.
\]
There is a unique feasible threshold $v_{SU}(c)$ whenever the initial SU-advantage region exists, and
\begin{equation}
\mathcal V_{SU}(c)
=
\begin{cases}
(v_{SU}(c),\bar v(c)),&0<c<c^*,\\
\varnothing,&c^*\le c<1/3,
\end{cases}
\label{eq:vsu-region}
\end{equation}
where $c^*\approx0.1529167823$ is the unique root in $(0,1/3)$ of
\[
17c^3+109c^2-89c+11=0.
\]
Positive network effects are therefore necessary for the initial SU advantage in this canonical microfoundation, not for private circumvention itself.

\subsection{Singleton-network specification}
\label{subsec:singleton-network-sensitivity}

The canonical network term assigns zero network benefit to a singleton compatibility group. R4 tests one pre-specified alternative specification in which every nonempty group, including a singleton, receives
\[
g_i=v\sum_{j\in G_i}q_j.
\]
No other primitive is changed.

Under this alternative, the one-way-adoption logic remains feasible at an exact witness, but the pre-adoption member gap satisfies
\begin{equation}
    \Phi_{S1}<0
\end{equation}
throughout the old canonical feasible domain. The political reversal route is therefore eliminated before private adoption occurs.

This is a specification-sensitivity result, not a universal theorem about singleton network benefits. R7 certifies only that the canonical zero-singleton convention is substantively relevant to the initial government ranking.

\subsection{Differentiated-demand portability test}
\label{subsec:differentiated-portability}

R5 tests one alternative demand environment fixed before solving it. The substitution parameter is set to $\gamma=1/2$, compatibility enters a quadratic utility system, and both quantity and price competition are solved under the same differentiated demand. This construction separates a change in strategic variable from a change in the demand microfoundation.

Let $X\in\{C,B\}$ index the differentiated Cournot and differentiated Bertrand versions. R5 proves throughout the frozen audit box
\begin{equation}
    \Phi^X
    =
    W_M^{SU,N,X}-W^{IS,X}<0,
    \label{eq:r5-negative-phi}
\end{equation}
and
\begin{equation}
    \mathscr E^X<0.
    \label{eq:r5-negative-e}
\end{equation}
Thus the member-market component that drives the canonical political effect changes sign in both versions. At the same time, outsider-only private adoption remains feasible and selection free at the common exact witness used in R5.

The failure therefore cannot be attributed to Bertrand competition alone. Nor does it establish that every differentiated-product model eliminates the mechanism. The certified conclusion is narrower: the political preference effect is not portable to this one pre-specified differentiated-demand environment even though one-way private adoption survives.

\subsection{Coalition-rule and symmetry dependence}
\label{subsec:institutional-scope}

Section~\ref{sec:coalition-stability} already shows that the high-$F$ symmetric stable set is sensitive to the blocking rule. Under strict blocking, all three SUs are stable because a common member is indifferent when switching partners. Under weak/Pareto blocking, the same indifference no longer prevents a deviation, and the high-$F$ stable set is empty.

A small market-size perturbation,
\[
m_1=m_2=1,
\qquad
m_3=1-\delta,
\]
breaks the equality:
\begin{equation}
    W_1^{12,N}-W_1^{13,N}
    =
    \delta(A-C)>0.
    \label{eq:robustness-asymmetry-gap}
\end{equation}
At the certified high-$F$ witness, the two-large-country union $\rho_{12}^{SU}$ is uniquely stable under both blocking concepts for every $0<\delta<1$. In the intermediate region, the unique-IS result is supported by strict payoff differences and survives the same perturbation over the certified range.

The institutional conclusion is therefore asymmetric. The intermediate full-bypass application is relatively robust within the canonical product-market model. The exact symmetric high-$F$ multiplicity is not.

\input{tables/generated/table_scope_robustness}

Figure~\ref{fig:assumption-dependence} summarizes this hierarchy.
\begin{figure}[t]
    \centering
    \includegraphics[width=\linewidth]{figures/generated/figure_04_assumption_dependence.pdf}
    \caption{Assumption dependence of the main results. One-way private adoption survives more of the specification checks than the initial SU advantage and the resulting government-ranking reversal. The high-$F$ stable-set multiplicity is additionally sensitive to the blocking rule and exact symmetry.}
    \label{fig:assumption-dependence}
\end{figure}

The robustness exercises support a deliberately limited interpretation. The firm-level circumvention logic is comparatively portable. The government-ranking effect requires the canonical market structure to generate a positive initial member-market component, and the exact stable-set correspondence also depends on the coalition institution. The paper therefore treats the success and failure boundaries as part of the result rather than as caveats to be hidden after the theorem.
